% Level-pq Lambert annihilator expository note --- pre-submission LaTeX source.
% Target venue: Mathematics Magazine (MAA). Style: amsart, amsplain bib.
% Compile: pdflatex main && bibtex main && pdflatex main && pdflatex main
% Or: latexmk -pdf main.tex

\documentclass[11pt,reqno]{amsart}

\usepackage[utf8]{inputenc}
\usepackage[T1]{fontenc}
\usepackage{amsmath,amsthm,amsfonts,amssymb}
\usepackage{mathrsfs}
\usepackage{microtype}
\usepackage[margin=1in]{geometry}
\usepackage{booktabs}
\usepackage{array}
\usepackage{xcolor}
\usepackage{hyperref}
\hypersetup{
    colorlinks=true,
    linkcolor=blue!50!black,
    citecolor=blue!50!black,
    urlcolor=blue!50!black
}
\usepackage{cleveref}

% Theorem-like environments
\newtheorem{theorem}{Theorem}[section]
\newtheorem{proposition}[theorem]{Proposition}
\newtheorem{lemma}[theorem]{Lemma}
\newtheorem{corollary}[theorem]{Corollary}

\theoremstyle{definition}
\newtheorem{remark}[theorem]{Remark}
\newtheorem{example}[theorem]{Example}

% Custom shortcuts
\newcommand{\Z}{\mathbb{Z}}
\newcommand{\R}{\mathbb{R}}
\newcommand{\Q}{\mathbb{Q}}
\newcommand{\HH}{\mathbb{H}}
\newcommand{\Eis}{\mathcal{E}}
\DeclareMathOperator{\Gal}{Gal}

\title{An explicit closed form for the four-divisor $E_2$ Lambert annihilator at
squarefree composite level: a note on autonomous rediscovery}

\author{Jouni Helminen}
\address{Independent researcher}
\email{jouni@helminen.co}

\date{May 25, 2026}

\subjclass[2020]{Primary 11F11, 11F20; Secondary 11M36, 68T20}

\keywords{Lambert series; Eisenstein series; quasi-modular form $E_2$;
level $\Gamma_0(N)$; Hardy--Ramanujan--Berndt identity; autonomous
mathematical rediscovery}

\begin{document}

\begin{abstract}
We give an explicit closed-form parameterization of the unique
(up to scale) linear relation among the four Lambert ``conjugate-pair
blocks''
\[
T_d(\alpha,\beta) := d\,\alpha\,L(d\alpha)
+ \frac{\beta}{d}\,L\!\left(\frac{\beta}{d}\right),
\qquad d \mid N,
\qquad L(x) := \sum_{n \ge 1}\frac{n}{e^{2nx}-1},
\]
on the hyperbola $\alpha\beta = \pi^2$ at squarefree composite level
$N = pq$ with distinct primes $p < q$. With
$g := \gcd(pq-1,\, q-p)$, $\delta := (q-p)/g$, and $m := (pq-1)/g$,
the integer-lowest-height coefficient vector is
\[
(c_1, c_p, c_q, c_{pq}) = (-\delta,\; m,\; -m,\; \delta).
\]
The derivation requires only the Hardy--Ramanujan--Berndt base identity
(1916) plus five lines of Gaussian elimination on a $3\times 4$ rational
system. The underlying structural fact---that exactly one linear
relation must exist---is a consequence of the classical dimension
formula $\dim \Eis_2(\Gamma_0(pq)) = 3$ for the Eisenstein subspace of
weight-$2$ modular forms at squarefree composite level. We verify the
formula numerically across sixteen instances to relative error
$\le 10^{-100}$ at \texttt{mpmath} precision $200$ digits, map
structural extensions (the formula is unique at $N = pq$, non-unique
at $N = pqr$ and $N = p^2 q$, and does not extend to $N = p^2$), and
offer the note as both a useful explicit parameterization and as an
instance of autonomous rediscovery of classical material by an
LLM-driven symbolic-search engine.
\end{abstract}

\maketitle

\section{Introduction}\label{sec:intro}

In their 1916 paper \emph{Asymptotic formulae in combinatory analysis}
\cite[eq.~(1.55)]{HardyRamanujan1916}, Hardy and Ramanujan recorded the
following identity, equivalent to the modular transformation law for
the weight-$2$ Eisenstein series $E_2$: for $\alpha,\beta > 0$ with
$\alpha\beta = \pi^2$,
\begin{equation}\label{eq:HRB}
\alpha\, L(\alpha) + \beta\, L(\beta)
= \frac{\alpha + \beta}{24} - \frac{1}{4},
\end{equation}
where $L(x) := \sum_{n \ge 1} n/(e^{2nx} - 1)$ is the standard Lambert
series. Berndt's commentary in \emph{Ramanujan's Notebooks Part~II}
\cite[Ch.~14, Entry~27(iii)]{BerndtNB2} traces this identity to
Ramanujan's second notebook and connects it directly to the $E_2$
inversion law
\begin{equation}\label{eq:E2}
E_2(-1/\tau) = \tau^2\, E_2(\tau) + \frac{6\tau}{\pi i},
\end{equation}
via the substitution $\tau = i\alpha/\pi$ and the standard
$q$-expansion $E_2(\tau) = 1 - 24\sum_{n \ge 1} n\, q^n/(1 - q^n)$
with $q = e^{2\pi i \tau}$.

A natural question is: given a level $N$, what is the dimension of the
linear-relation space spanned by the conjugate-pair blocks
\begin{equation}\label{eq:Td}
T_d(\alpha,\beta) := d\,\alpha\,L(d\alpha)
+ \frac{\beta}{d}\,L\!\left(\frac{\beta}{d}\right)
\end{equation}
indexed by $d \mid N$? For $N = 1$ there is one block and
\cref{eq:HRB} determines its value but gives no annihilator. For
$N = p$ prime, there are two blocks $\{T_1, T_p\}$ but
\cref{eq:HRB} applied at the two conjugate pairs $(\alpha,\beta)$ and
$(p\alpha, \beta/p)$ produces no non-trivial linear relation (both
blocks are affine in $(\alpha,\beta)$ with linearly independent
coefficient vectors).

The first non-trivial case is therefore $N = pq$ with $p \ne q$ both
prime: four blocks indexed by $d \in \{1, p, q, pq\}$. The main result
of this note is that this case admits exactly one non-trivial linear
relation, with explicit closed-form coefficients.

\begin{theorem}\label{thm:main}
Let $p < q$ be distinct primes. With $g := \gcd(pq-1,\, q-p)$,
$\delta := (q-p)/g$, and $m := (pq-1)/g$, the unique (up to scale)
integer-lowest-height annihilator of $\{T_1, T_p, T_q, T_{pq}\}$ on the
hyperbola $\alpha\beta = \pi^2$ is
\begin{equation}\label{eq:main}
-\delta\, T_1 + m\, T_p - m\, T_q + \delta\, T_{pq} = 0.
\end{equation}
\end{theorem}

The remainder of this note is an exposition of \cref{eq:main}: we
derive it in five lines (\cref{sec:derivation}), connect it to the
classical $E_2$ Eisenstein-subspace dimension formula
(\cref{sec:modular}), tabulate sixteen numerically verified instances
(\cref{sec:numerical}), map its structural extensions to richer
divisor sets (\cref{sec:extensions}), and close with a provenance note
on how the formula was rediscovered (\cref{sec:provenance}).

The underlying mathematics is entirely classical. The fact that
exactly one such relation must exist at $N = pq$ follows immediately
from $\dim \Eis_2(\Gamma_0(pq)) = 3$. The explicit
$\gcd(pq-1, q-p)$ parameterization for the integer-lowest-height
representative may not appear in this form in the standard
references; if it does, we have not located it. The note is offered
both as a convenient explicit form and as a worked example of
LLM-assisted autonomous rediscovery of a classical structural
identity.

\section{Closed-form coefficients}\label{sec:derivation}

\subsection{The five-line derivation}\label{sec:fiveline}

Identity~\eqref{eq:HRB} collapses each conjugate-pair block to a
three-parameter affine function of $(\alpha,\beta)$ alone.
Substituting the conjugate pair $(d\alpha,\, \beta/d)$ into
\cref{eq:HRB}---which is permissible because
$d\alpha \cdot \beta/d = \alpha\beta = \pi^2$---we obtain
\begin{equation}\label{eq:Td-affine}
T_d(\alpha,\beta) = \frac{d\alpha + \beta/d}{24} - \frac{1}{4}
= \frac{d\,\alpha}{24} + \frac{\beta}{24\,d} - \frac{1}{4}.
\end{equation}

A linear combination $\sum_{d \mid pq} c_d\, T_d$ therefore equals
\begin{equation}\label{eq:combination}
\frac{\alpha}{24}\sum_{d} d\,c_d
+ \frac{\beta}{24}\sum_{d} \frac{c_d}{d}
- \frac{1}{4}\sum_d c_d.
\end{equation}

The combination vanishes identically in $(\alpha,\beta)$ iff the three
coefficient sums all vanish:
\begin{equation}\label{eq:three-constraints}
\sum_d c_d = 0,\qquad
\sum_d d\,c_d = 0,\qquad
\sum_d \frac{c_d}{d} = 0.
\end{equation}

With $d \in \{1, p, q, pq\}$ and unknowns $(c_1, c_p, c_q, c_{pq})$,
labelling \cref{eq:three-constraints} as $E_1$, $E_2$, $E_3$
respectively, this is a $3 \times 4$ rational linear system. Five
lines of Gaussian elimination solve it:
\begin{enumerate}
\item $E_2 - E_1$:\quad
$(p-1)c_p + (q-1)c_q + (pq-1)c_{pq} = 0$.
\item $-pq(E_3 - E_1)$:\quad
$q(p-1)c_p + p(q-1)c_q + (pq-1)c_{pq} = 0$.
\item Step~1 $-$ Step~2:\quad
$-(p-1)(q-1)(c_p + c_q) = 0$, so $c_q = -c_p$.
\item Substituting into Step~1:\quad
$(pq-1)c_{pq} = (q-p)\,c_p$, i.e.\
$m\,(q-p) = \delta\,(pq-1)$ with $c_p =: m$ and $c_{pq} =: \delta$.
\item Substituting into $E_1$ with $c_q = -c_p$:\quad
$c_1 = -c_{pq} = -\delta$.
\end{enumerate}

The integer-lowest-height representative is then obtained by dividing
through by $g := \gcd(pq-1,\, q-p)$, yielding $\delta = (q-p)/g$,
$m = (pq-1)/g$, and the coefficient vector
$(c_1, c_p, c_q, c_{pq}) = (-\delta,\, m,\, -m,\, \delta)$ as
claimed. This completes the proof of \cref{thm:main}.\qed

\subsection{Worked example: $N = 33 = 3 \cdot 11$}

Here $p = 3$, $q = 11$, so $q - p = 8$, $pq - 1 = 32$, and
$g = \gcd(32, 8) = 8$. Therefore $\delta = 8/8 = 1$ and $m = 32/8 = 4$,
giving
\begin{equation}\label{eq:N33}
-T_1(\alpha,\beta) + 4\, T_3(\alpha,\beta)
- 4\, T_{11}(\alpha,\beta) + T_{33}(\alpha,\beta) = 0.
\end{equation}
Expanding the blocks explicitly,
\begin{multline*}
\alpha\bigl[-L(\alpha) + 12\,L(3\alpha) - 44\,L(11\alpha)
+ 33\,L(33\alpha)\bigr]\\
+ \beta\bigl[-L(\beta) + \tfrac{4}{3}\,L(\beta/3)
- \tfrac{4}{11}\,L(\beta/11) + \tfrac{1}{33}\,L(\beta/33)\bigr] = 0.
\end{multline*}
Numerically, at the canonical cell $\alpha = 1$, $\beta = \pi^2$, the
left-hand side evaluates to magnitude $\le 10^{-199}$ at
\texttt{mpmath} working precision $200$ digits, against blocks of
magnitude $\sim 10^{-1}$---relative error $\le 10^{-199}$.

\section{Connection to classical modular forms}\label{sec:modular}

\subsection{The dimension formula}

For squarefree $N = q_1 q_2 \cdots q_t$ with $t$ distinct prime
factors, the Eisenstein subspace of weight-$2$ modular forms on
$\Gamma_0(N)$ has dimension
\begin{equation}\label{eq:dim}
\dim \Eis_2(\Gamma_0(N)) = 2^t - 1.
\end{equation}
This is a standard fact; see Diamond--Shurman \cite[\S3.5]{DiamondShurman}
or Stein \cite[Ch.~5]{Stein}. For $t = 2$ (the squarefree composite
case $N = pq$) this gives $\dim \Eis_2(\Gamma_0(pq)) = 3$.

The four functions $\{E_2(\tau),\, E_2(p\tau),\, E_2(q\tau),\,
E_2(pq\tau)\}$ are weight-$2$ Eisenstein-style forms on
$\Gamma_0(pq)$. Each $E_2(d\tau)$ is a quasi-modular form on
$\Gamma_0(d) \supseteq \Gamma_0(pq)$, and the standard old-form
construction places it in the Eisenstein subspace of
$M_2(\Gamma_0(pq))$ after subtracting the appropriate quasi-modular
anomaly; see \cite[\S5.6]{DiamondShurman}. Because four candidate
forms live in a three-dimensional space, they must satisfy at least
one linear relation. The argument of \cref{sec:fiveline}---which used
only \cref{eq:HRB}, not the Eisenstein-subspace framework---shows
that the linear-relation space is exactly one-dimensional.

The LMFDB database \cite{LMFDB} confirms the dimension count at
$N = 33$: the space $M_2(\Gamma_0(33))$ has total dimension~$9$,
cuspidal subspace dimension~$6$, and Eisenstein subspace dimension~$3$.

\subsection{The Hardy--Ramanujan--Berndt identity in this light}

Identity~\eqref{eq:HRB} is precisely the manifestation, at the level
of one conjugate pair of Lambert blocks, of the transformation
law~\eqref{eq:E2} for $E_2$. The anomaly $6\tau/(\pi i)$ in
\cref{eq:E2} becomes the constant correction $-1/4$ and the linear
piece $(\alpha + \beta)/24$ in \cref{eq:HRB}. The fact that this
anomaly is \emph{itself a polynomial of low degree in
$(\alpha, \beta)$}---specifically of degree at most one---is what
reduces the four-divisor annihilator equation to only three linear
constraints (\cref{eq:three-constraints}), and hence to a kernel of
dimension $4 - 3 = 1$ in the generic case.

Said differently: the Hardy--Ramanujan--Berndt identity is the
dimension reduction $\dim \Eis_2(\Gamma_0(pq)) = 3$ in disguise,
working at the level of Lambert summands rather than modular forms.

\section{Numerical verification}\label{sec:numerical}

We verified \cref{eq:main} across sixteen instances using
\texttt{mpmath} at working precision $200$ digits, evaluating the
left-hand side at five cells on the hyperbola $\alpha\beta = \pi^2$
chosen to stress the formula at small $\alpha$, large $\alpha$, and
the symmetric point $\alpha = \beta = \pi$. The strict acceptance bar
is relative error $\le 10^{-100}$, where the relative error is
measured against $\max_d |c_d \cdot T_d|$ (the cancellation scale).

\Cref{tab:verification} summarizes seven main instances plus nine
adversarial edge cases.

\begin{table}[h]
\centering
\small
\setlength{\tabcolsep}{5pt}
\begin{tabular}{rcclr}
\toprule
$N$ & $(p, q)$ & Coefficients $(c_1, c_p, c_q, c_{pq})$
& $\,g,\, \delta,\, m\,$ & Worst rel.\,err.\ / 5 cells \\
\midrule
$33$  & $(3, 11)$  & $(-1, 4, -4, 1)$       & $8, 1, 4$    & $\le 10^{-199}$ \\
$35$  & $(5, 7)$   & $(-1, 17, -17, 1)$     & $2, 1, 17$   & $2.76\times 10^{-200}$ \\
$51$  & $(3, 17)$  & $(-7, 25, -25, 7)$     & $2, 7, 25$   & $2.85\times 10^{-200}$ \\
$65$  & $(5, 13)$  & $(-1, 8, -8, 1)$       & $8, 1, 8$    & $4.42\times 10^{-200}$ \\
$77$  & $(7, 11)$  & $(-1, 19, -19, 1)$     & $4, 1, 19$   & $1.27\times 10^{-200}$ \\
$87$  & $(3, 29)$  & $(-13, 43, -43, 13)$   & $2, 13, 43$  & $1.99\times 10^{-199}$ \\
$95$  & $(5, 19)$  & $(-7, 47, -47, 7)$     & $2, 7, 47$   & $4.88\times 10^{-200}$ \\
$115$ & $(5, 23)$  & $(-3, 19, -19, 3)$     & $6, 3, 19$   & $1.07\times 10^{-199}$ \\
$143$ & $(11, 13)$ & $(-1, 71, -71, 1)$     & $2, 1, 71$   & $1.93\times 10^{-200}$ \\
$323$ & $(17, 19)$ & $(-1, 161, -161, 1)$   & $2, 1, 161$  & $6.99\times 10^{-134}$ \\
$187$ & $(11, 17)$ & $(-1, 31, -31, 1)$     & $6, 1, 31$   & $9.60\times 10^{-200}$ \\
$391$ & $(17, 23)$ & $(-1, 65, -65, 1)$     & $6, 1, 65$   & $3.46\times 10^{-110}$ \\
$6$   & $(2, 3)$   & $(-1, 5, -5, 1)$       & $1, 1, 5$    & $4.02\times 10^{-200}$ \\
$10$  & $(2, 5)$   & $(-1, 3, -3, 1)$       & $3, 1, 3$    & $1.91\times 10^{-200}$ \\
$14$  & $(2, 7)$   & $(-5, 13, -13, 5)$     & $1, 5, 13$   & $2.70\times 10^{-200}$ \\
$22$  & $(2, 11)$  & $(-3, 7, -7, 3)$       & $3, 3, 7$    & $1.22\times 10^{-199}$ \\
\bottomrule
\end{tabular}
\medskip
\caption{Numerical verification of \cref{eq:main} across sixteen
$N = pq$ instances at \texttt{mpmath} working precision $200$ digits.
Worst relative error is the maximum, across five test cells on
$\alpha\beta = \pi^2$, of $|{\rm LHS}|/\max_d|c_d\,T_d|$.}
\label{tab:verification}
\end{table}

The bar-marginal entries at $N = 323$ (twin primes $17 \cdot 19$) and
$N = 391$ (far-apart primes $17 \cdot 23$) reflect numerical
cancellation eating $\sim 66$ and $\sim 90$ digits respectively in
the most extreme cell; both remain comfortably inside the $10^{-100}$
bar. All other entries achieve relative error $\le 10^{-199}$, i.e.\
essentially the working-precision floor.

The formula is symmetric under $p \leftrightarrow q$: the substitution
flips $\delta \to -\delta$ and the labels $c_p \leftrightarrow c_q$,
leaving the relation invariant up to an overall sign. The derivation
in \cref{sec:fiveline} nowhere assumed $p > 2$; the small-prime
entries $N \in \{6, 10, 14, 22\}$ confirm this.

\section{Structural extensions}\label{sec:extensions}

What happens at richer divisor sets? \Cref{sec:fiveline} shows that
any divisor set $D = \{d_1, \ldots, d_n\}$ on which $T_d$ is defined
yields the same three linear constraints \eqref{eq:three-constraints},
whose kernel has dimension $n - \mathrm{rank}(M)$ where $M$ is the
$3 \times n$ matrix with rows $(1, \ldots, 1)$,
$(d_1, \ldots, d_n)$, $(1/d_1, \ldots, 1/d_n)$. The interesting
structural extensions are:

\paragraph{Squarefree $N = pqr$ (three distinct primes).}
Eight divisors $\{1, p, q, r, pq, pr, qr, pqr\}$, a $3 \times 8$
rational system, kernel dimension~$5$. The closed-form
parameterization of \cref{thm:main} does \emph{not} extend uniquely;
the kernel is non-unique. Numerically verified at
$N \in \{30,\, 105,\, 165\}$. The first basis vector is typically a
``sub-divisor annihilator'' supported on $\{1, p, q, r\}$ alone; the
lifted \cref{thm:main} annihilators on the embedded $pq$-, $pr$-,
$qr$-sublattices also live inside the kernel.

\paragraph{Prime power $N = p^2$.}
Three divisors $\{1, p, p^2\}$. The $3 \times 3$ constraint matrix is
Vandermonde-type in $(1, p, p^2)$ and $(1, 1/p, 1/p^2)$, of full
rank~$3$, kernel dimension~$0$. \emph{The family does not extend to
prime powers.} Verified at $N \in \{49, 121\}$. The construction
structurally requires at least two distinct prime factors.

\paragraph{Mixed $N = p^2 q$.}
Six divisors $\{1, p, q, p^2, pq, p^2 q\}$, kernel dimension~$3$. The
closed form is non-unique. Verified at $N \in \{18, 45\}$. As at
$pqr$, sub-divisor annihilators and lifts of the \cref{thm:main}
$p \cdot q$ annihilator both populate the kernel.

The clean uniqueness of \cref{thm:main} is therefore special to the
squarefree composite case $N = pq$.

\begin{remark}\label{rem:caveat}
A subtler point: at $N = pqr$ the Eisenstein-subspace dimension
formula \eqref{eq:dim} gives $\dim \Eis_2(\Gamma_0(pqr)) = 2^3 - 1 = 7$,
so the eight forms $\{E_2(d\tau) : d \mid pqr\}$ satisfy exactly one
linear relation at the level of modular forms---yet our Lambert-block
construction yields a kernel of dimension~$5$, not~$1$. The
discrepancy is because \cref{eq:HRB} captures only the affine
truncation of the full $E_2$ $q$-expansion under modular inversion;
it loses information about the higher-order terms that modular-form
linear-relation theory tracks in full. At $N = pq$ the kernel
dimensions coincide ($1 = 1$); at higher $N$ they generically
diverge. This is a useful caveat: the Hardy--Ramanujan--Berndt
identity is a \emph{strictly weaker} constraint than full
$E_2$-modularity, and the Lambert-annihilator picture and the
Eisenstein-subspace linear-relation picture only agree at $N = pq$.
\end{remark}

\section{Provenance: autonomous rediscovery by symbolic search}\label{sec:provenance}

The closed-form formula~\eqref{eq:main} was rediscovered autonomously
by Project~C, an LLM-driven symbolic-search engine for
Ramanujan-style identities developed during the 2026
\emph{Ramanujan mining} research run \cite{ProjectC}. The engine's
basis of operators---\verb|mellin_inversion|, \verb|master_theorem|,
\verb|modular_transform_with_correction|, and others---was applied
iteratively to a small library of seed identities under an
\texttt{mpmath} cost model with a strict $\le 10^{-30}$ relative-error
acceptance bar. The Level-$33$ instance was the first the engine
surfaced (after several thousand seed-operator compositions); a
prompted exploration cycle then enumerated the seven additional
instances $N \in \{35, 51, 65, 77, 87, 95, 115\}$ from a
hand-tabulated guess of the integer-lowest-height coefficient
pattern, yielding the explicit $\gcd(pq-1, q-p)$ form.

Crucially, the engine did \emph{not} use the classical dimension
formula~\eqref{eq:dim} or the Eisenstein-subspace framework of
\cref{sec:modular}. It worked entirely from \cref{eq:HRB} and
elementary linear algebra over the rationals. The structural
connection to $\dim \Eis_2(\Gamma_0(pq)) = 3$ was identified only
\emph{after} the formula was produced, by an adversarial verification
subagent performing a literature audit. Where helpful, the run also
made use of OpenAI's frontier-model interface as a complementary
sounding board \cite{ChatGPTPro}.

A subsequent literature diff \cite[diary~58]{ProjectC} across six
volumes of Berndt's \emph{Ramanujan's Notebooks} and \emph{Lost
Notebook} commentary \cite{BerndtNB2, AndrewsBerndt}, Berndt--Yee
\cite{BerndtYee}, Berndt--Dixit \cite{BerndtDixit}, plus the LMFDB
database, the OEIS, and approximately twelve targeted arXiv and AMS
searches, did not locate the explicit $\gcd(pq-1, q-p)$
parameterization in this exact form. The underlying \emph{structure}---a
one-dimensional linear-relation space among the four Eisenstein
candidates at squarefree composite level---is, however, classical
and well-documented.

We therefore offer this note in two registers. First, as a useful
explicit closed form for a relation that is \emph{guaranteed to
exist} by Eisenstein-subspace dimension counting but is rarely
written down in this elementary form, and which evaluates exactly to
zero at every $(\alpha, \beta)$ on $\alpha\beta = \pi^2$. Second, as
a worked instance of autonomous rediscovery: an LLM-driven engine,
operating with no prior modular-forms machinery and only a small
basis of structural operators, rediscovered in explicit form a
classical structural identity that has been implicit in the
modular-forms literature for decades. We regard the latter as the
more interesting contribution. The mathematics is classical; what is
novel is the demonstration that a small, structured operator basis
plus a strict numerical acceptance gate suffices to surface this
structure autonomously.

The full computational verification---numerical code, raw results,
and the literature audit---is available at \cite{ProjectC} in
\verb|analysis/30_level33_lambert_e2_annihilator_verification.py|
and \verb|analysis/32_level_pq_family_adversarial_verification.py|.

\section*{Acknowledgments}

The author thanks the developers of \texttt{mpmath}, \texttt{sympy},
the LMFDB project, and the maintainers of the open-access portions
of Berndt's \emph{Ramanujan's Notebooks} series.

\bibliographystyle{amsplain}
\bibliography{refs}

\end{document}
